\documentclass[11pt, a4paper]{article}

% --- Encodage et Langue ---
\usepackage[utf8]{inputenc}
\usepackage[T1]{fontenc}
\usepackage[french]{babel} % Adapte la typographie (guillemets, espaces)

% --- Mise en page ---
\usepackage[margin=2.5cm]{geometry} % Marges équilibrées

% --- Mathématiques ---
\usepackage{amsmath, amssymb, amsthm}


\title{Note de Synthèse}
\date{}

\begin{document}
\maketitle

\section{Dixon \& Coles}
Dans cet article \cite{dixon} Dixon et Coles se proposent d'étudier un modèle de 
prédiction permettant de prédire avec précision les probabilités 
de l'issue des matchs de football. 
Pour décrire le modèle, il faut ce baser sur le modèle de Maher (1982) 
qui suppose que les buts marqués par les équipes à domicile et à 
l'extérieur suivent des lois de Poisson indépendantes. 
D'autres modèles ont pu être proposés; Reep et al(1971) a proposé à l'instar de 
Moroney(1956) un modèle basé sur la binomiale négative. 
Cependant cette approche a été abandonnée par la suite ne trouvant 
aucun moyen de prédire les résultats.

Le modèle de Maher utilise donc une loi Poisson dont la moyenne est 
modélisée en fonction des performances présentes des équipes respectives.
Afin d'améliorer et de tester ce modèle, de nombreux papiers traite de différents 
paramètres, l'exclusion d'un joueur, le fait de jouer à domicile ou à l'extérieur, etc.


\subsection{Données}

L'étude s'appuie sur 6 629 résultats de matchs anglais (championnat et coupes) 
entre 1992 et 1995. Contrairement à d'autres approches plus complexes, les auteurs 
font le choix délibéré de n'utiliser que les scores finaux, ignorant volontairement 
des données comme l'identité des buteurs ou le moment des buts pour garantir 
la robustesse statistique du modèle.

Les données disponibles, qui comprennent $6 629$ résultats de matchs de championnat 
et de coupe (temps plein) des saisons $1992-1993$, $1993-1994$ et $1994-1995$, incluent 
pour chaque match le score à domicile et le score à l'extérieur. Les données de 
$1995-1996$ sont également disponibles, mais servent d'échantillon de validation 
pour tester l'utilité du modèle lorsqu'il est utilisé comme base pour une stratégie 
de paris.

\subsection{Modèle de Dixon et Coles}

Dans un premier temps, les auteurs mènenent une analyse empirique des données, 
en se concentrant sur les résultats à domicile, à l'extérieur et nuls. Ils constatent que 
l'ajustement d'une distribution de Poisson aux scores agrégés à domicile et à 
l'extérieur révèle que, quel que soit le critère, le modèle de Poisson correspond 
presque parfaitement. 

\subsection{Détail du Modèle et inférences}

Afin de développer au mieux leur modèle, les auteurs précisent certains éléments 
importants :
\begin{itemize}
    \item différentes capacités des deux équipes lors d'un match.
    \item Il convient de tenir compte du fait que les équipes jouant à domicile 
    bénéficient généralement d'un certain avantage, le fameux « effet terrain ».
    \item pour estimer la capacité d'une équipe, il est logique de se baser sur 
    les performances récentes de l'équipe
    \item Pour mieux détailler les capacités d'une équipe, il est nécessaire de les différencier.
    L'attaque et la défense.
    \item Pour résumer les performances d'une équipe à partir de ses résultats récents, 
    il convient de tenir compte du niveau des équipes qu'elle a affrontées.
\end{itemize}
Ils se basent donc sur le modèle de Maher(1982) en apportant certains ajustement afin 
de pouvoir traiter des données incomplètes, de pouvoir tenir compte de l'évolution 
des performances des équipes.


\section{Leonardo Egidi et Nicola Torelli}


\section{Carloni et al.}


\newpage
\nocite{*}
\bibliographystyle{plain}
\bibliography{references}
\end{document}